% =====================================================================
%  CHAPITRE 12 — VOCABULAIRE : STRATÉGIE ET THÈMES
% =====================================================================
\chapter{Vocabulaire : apprendre les mots qui servent}

\section{La stratégie : 3\,000 mots suffisent pour tout dire}
Les études de fréquence sont formelles :

\begin{tableau}[Couverture d'un texte ordinaire selon la taille du vocabulaire]
\begin{tabularx}{\linewidth}{@{}p{3.6cm}p{3.0cm}X@{}}
\entete{Mots connus} & \entete{Couverture} & \entete{Ce que vous pouvez faire} \tabularnewline
\hline
100 mots les plus fréquents & 50\,\% & Presque rien, mais c'est la structure de toute phrase. \tabularnewline
500 mots & 75\,\% & Phrases de survie, se débrouiller dans un magasin. \tabularnewline
1\,000 mots & 85\,\% & Conversation quotidienne simple. \tabularnewline
2\,000 mots & 90\,\% & Comprendre un journal, suivre une série. \tabularnewline
\textbf{3\,000 mots} & \textbf{95\,\%} & \textbf{Conversation courante complète. C'est l'objectif du mois 9.} \tabularnewline
5\,000 mots & 98\,\% & Lire un roman, travailler en anglais. \tabularnewline
20\,000 mots & 99\,\% & Niveau d'un natif cultivé --- inutile à viser. \tabularnewline
\end{tabularx}
\end{tableau}

\begin{astuce}[N'apprenez pas « des mots », apprenez des blocs]
\en{decision} seul ne sert à rien. \en{make a decision} sert immédiatement.
Apprenez toujours le mot \textbf{avec son compagnon} :\\
\en{make a decision, make a mistake, make progress, take a photo, take a break,
do your best, do the shopping, have a look, have a rest, get married, get lost,
pay attention, keep in touch, catch a cold, spend time, feel like doing}.
\end{astuce}

\section{Votre avantage secret : les mots que vous connaissez déjà}
Entre 30 et 45\,\% du vocabulaire anglais vient du français (souvent via le normand après
1066). Voici 120 mots anglais que vous connaissez \textbf{déjà}, à condition de les
prononcer à l'anglaise.

\begin{tableau}[120 mots transparents --- lisez-les à voix haute]
\footnotesize
\begin{multicols}{3}
\begin{anglais}
\begin{itemize}\setlength{\itemsep}{1pt}
\item able --- capable
\item accept --- accepter
\item accident --- accident
\item action --- action
\item active --- actif
\item address --- adresse
\item admirable
\item adventure --- aventure
\item advice --- conseil
\item age --- âge
\item agree --- être d'accord
\item almost --- presque
\item already --- déjà
\item animal --- animal
\item announce --- annoncer
\item answer --- réponse
\item apart --- à part
\item apple --- pomme (piège !)
\item argument --- argument
\item arrive --- arriver
\item article --- article
\item attention --- attention
\item attitude
\item automatic
\item available --- disponible
\item average --- moyen
\item avoid --- éviter
\item balance --- équilibre
\item bank --- banque
\item base --- base
\item beautiful --- beau
\item because --- parce que
\item begin --- commencer
\item believe --- croire
\item benefit --- bénéfice
\item bicycle --- bicyclette
\item brave --- courageux
\item cable
\item calendar --- calendrier
\item calm --- calme
\item camera --- appareil photo
\item capital --- capitale
\item careful --- prudent
\item central --- central
\item certain --- certain
\item chance --- chance
\item change --- changer
\item chapter --- chapitre
\item charity --- charité
\item circle --- cercle
\item city --- cité, ville
\item civil --- civil
\item clear --- clair
\item climate --- climat
\item collection
\item colour / color
\item comfortable
\item common --- commun
\item company --- entreprise
\item compare --- comparer
\item complete --- complet
\item complicated
\item condition
\item confident
\item consequence
\item contact
\item content --- contenu
\item continue --- continuer
\item control --- contrôler
\item conversation
\item copy --- copie
\item correct --- correct
\item country --- pays
\item courage
\item create --- créer
\item crisis --- crise
\item critic --- critique
\item culture
\item curious --- curieux
\item dance --- danse
\item danger
\item decide --- décider
\item declare --- déclarer
\item decorate --- décorer
\item defence / defense
\item delicate
\item demand --- demande
\item democracy
\item demonstrate
\item department
\item depend --- dépendre
\item describe --- décrire
\item desert --- désert
\item desire --- désir
\item detail
\item develop
\item dictionary
\item difference
\item difficult
\item digest
\item direct --- direct
\item disappear
\item discover
\item discuss
\item disease --- maladie (piège !)
\item distance
\item distribute
\item document
\item domestic
\item dominant
\item donate
\item double
\item dramatic
\item economy
\item effect --- effet
\item efficient
\item effort
\item elegant
\item element
\item emotion
\item employ --- employer
\item empty --- vide (piège !)
\item energy
\item enormous
\item enter --- entrer
\item environment
\item equal --- égal
\item essential
\item establish
\item estimate
\item exact --- exact
\item examine
\item excellent
\item except --- sauf
\item excuse --- excuse
\item exercise --- exercice
\item exist
\item experience
\item explain
\item express
\item extraordinary
\item extreme
\item fabric --- tissu (piège !)
\item face --- visage
\item fact --- fait
\item factor
\item family
\item famous --- célèbre
\item favourite
\item festival
\item figure
\item final --- final
\item finance
\item firm --- entreprise
\item flexible
\item fluent
\item focus
\item force --- force
\item forest --- forêt
\item forget --- oublier (piège !)
\item formal --- formel
\item fortune
\item fragile
\item free --- libre / gratuit
\item frequent
\item fresh --- frais
\item front --- devant
\item fruit --- fruit
\item function
\item fundamental
\item furniture --- meubles
\item future --- futur
\item gallery
\item garden --- jardin
\item general
\item generous
\item gentle --- doux
\item gesture --- geste
\item global
\item government
\item grateful --- reconnaissant
\item group --- groupe
\item guarantee
\item guide
\item habit --- habitude
\item harmonious
\item honest --- honnête
\item honour --- honneur
\item horrible
\item hospital
\item hotel
\item humour
\item idea --- idée
\item identical
\item ignore --- ignorer
\item illegal
\item illustrate
\item imagine
\item immediate
\item immense
\item immigrant
\item immune
\item impact
\item impatient
\item important
\item impossible
\item impression
\item improve --- améliorer
\item include
\item increase
\item incredible
\item independent
\item individual
\item industrial
\item influence
\item inform
\item information
\item ingredient
\item inhabit --- habiter (piège !)
\item initial
\item innocent
\item insist
\item inspiration
\item install
\item instant
\item instinct
\item institution
\item instruct
\item instrument
\item insurance --- assurance
\item intelligent
\item intense
\item intention
\item interest
\item international
\item interpret
\item interrupt
\item interview
\item intimate
\item introduce
\item invent
\item invest
\item invisible
\item invitation
\item irregular
\item isolate
\item issue --- problème
\item January --- janvier
\item journey --- voyage
\item joy --- joie
\item judge --- juge
\item juice --- jus
\item July --- juillet
\item June --- juin
\item junior
\item justice
\item justify
\item keen --- passionné (piège !)
\end{itemize}
\end{anglais}
\end{multicols}
\end{tableau}

\begin{piege}[Les transparents trompeurs]
Dans cette liste, certains mots sont des \textbf{faux amis} partiels : \en{disease}
(maladie, pas « désaise »), \en{empty} (vide), \en{fabric} (tissu), \en{forget} (oublier),
\en{inhabit} (habiter un lieu, pas « habiter » au sens de vivre : on dit \en{live in}),
\en{keen} (passionné), \en{free} (libre \emph{et} gratuit), \en{issue} (problème, pas
« issue » de secours --- \en{exit}). Le chapitre 13 liste les 150 faux amis complets.
\end{piege}

\section{Les familles de mots : un mot appris = cinq mots gagnés}
\begin{tableau}[Exemples de familles]
\footnotesize
\begin{longtable}{@{}p{2.6cm}p{3.2cm}p{3.2cm}p{3.0cm}p{2.0cm}@{}}
\entete{Verbe} & \entete{Nom (personne)} & \entete{Nom (chose/action)} & \entete{Adjectif} & \entete{Adverbe} \tabularnewline
\endfirsthead
\entete{Verbe} & \entete{Nom (pers.)} & \entete{Nom (chose)} & \entete{Adjectif} & \entete{Adverbe} \tabularnewline
\endhead
\hline
\en{decide} & \en{decision-maker} & \en{decision} & \en{decisive, decided} & \en{decisively} \tabularnewline
\en{develop} & \en{developer} & \en{development} & \en{developed, developing} & --- \tabularnewline
\en{employ} & \en{employer, employee} & \en{employment} & \en{employable, unemployed} & --- \tabularnewline
\en{explain} & --- & \en{explanation} & \en{explanatory} & --- \tabularnewline
\en{govern} & \en{governor} & \en{government} & \en{governmental} & --- \tabularnewline
\en{inform} & --- & \en{information} & \en{informative, informed} & --- \tabularnewline
\en{manage} & \en{manager} & \en{management} & \en{manageable} & --- \tabularnewline
\en{organise} & \en{organiser} & \en{organisation} & \en{organised} & --- \tabularnewline
\en{produce} & \en{producer} & \en{production, product} & \en{productive} & \en{productively} \tabularnewline
\en{succeed} & \en{successor} & \en{success} & \en{successful} & \en{successfully} \tabularnewline
\en{teach} & \en{teacher} & \en{teaching} & --- & --- \tabularnewline
\en{understand} & --- & \en{understanding} & \en{understandable} & \en{understandably} \tabularnewline
\end{longtable}
\end{tableau}

\begin{tableau}[Préfixes utiles (ils doublent votre vocabulaire)]
\begin{tabularx}{\linewidth}{@{}p{2.6cm}p{3.0cm}X@{}}
\entete{Préfixe} & \entete{Sens} & \entete{Exemples} \tabularnewline
\hline
\en{un-} & contraire & \en{unhappy, unsure, unfair, unusual, unemployed} \tabularnewline
\en{in- / im- / il- / ir-} & contraire & \en{impossible, incorrect, illegal, irregular, informal} \tabularnewline
\en{dis-} & contraire & \en{disagree, disappear, dislike, dishonest} \tabularnewline
\en{mis-} & mal & \en{mistake, misunderstand, misspell, mislead} \tabularnewline
\en{re-} & à nouveau & \en{rewrite, rebuild, redo, return, recycle} \tabularnewline
\en{pre-} & avant & \en{prepare, prevent, predict, preview} \tabularnewline
\en{over-} & trop & \en{overwork, overeat, overestimate, overtime} \tabularnewline
\en{under-} & pas assez, dessous & \en{underpaid, underestimate, underwear, underground} \tabularnewline
\en{out-} & dépasser & \en{outrun, outnumber, output} \tabularnewline
\en{co-} & ensemble & \en{cooperate, coexist, co-author} \tabularnewline
\en{self-} & soi-même & \en{self-confidence, self-employed, selfish} \tabularnewline
\en{anti-} & contre & \en{antisocial, antibiotic} \tabularnewline
\en{non-} & non & \en{non-stop, nonsense, non-smoker} \tabularnewline
\end{tabularx}
\end{tableau}

\begin{tableau}[Suffixes utiles]
\begin{tabularx}{\linewidth}{@{}p{2.8cm}p{3.4cm}X@{}}
\entete{Suffixe} & \entete{Crée} & \entete{Exemples} \tabularnewline
\hline
\en{-er / -or} & personne qui fait & \en{teacher, worker, actor, doctor, writer} \tabularnewline
\en{-ist} & spécialiste & \en{scientist, artist, dentist, economist} \tabularnewline
\en{-ian} & métier, origine & \en{musician, politician, Nigerian, Ghanaian} \tabularnewline
\en{-ee} & personne qui reçoit & \en{employee, trainee, refugee} \tabularnewline
\en{-ment} & nom d'action & \en{development, agreement, movement} \tabularnewline
\en{-tion / -sion} & nom d'action & \en{education, decision, information} \tabularnewline
\en{-ness} & nom d'état & \en{happiness, kindness, darkness, illness} \tabularnewline
\en{-ity} & nom de qualité & \en{ability, possibility, university} \tabularnewline
\en{-ance / -ence} & nom d'état & \en{importance, difference, patience} \tabularnewline
\en{-ful} & plein de & \en{beautiful, useful, careful, helpful} \tabularnewline
\en{-less} & sans & \en{useless, careless, homeless, wireless} \tabularnewline
\en{-able / -ible} & qui peut être & \en{possible, comfortable, available, reliable} \tabularnewline
\en{-ous} & qui a la qualité & \en{dangerous, famous, serious, generous} \tabularnewline
\en{-y} & caractérisé par & \en{rainy, sunny, dirty, happy, tasty} \tabularnewline
\en{-al} & relatif à & \en{national, personal, natural, musical} \tabularnewline
\en{-ish} & approximatif, nationalité & \en{reddish, thirtyish, British, childish} \tabularnewline
\en{-ly} & adverbe & \en{quickly, carefully, usually} \tabularnewline
\end{tabularx}
\end{tableau}

\section{Vocabulaire thématique}
Vingt thèmes, les mots dont vous aurez réellement besoin. Chaque entrée donne la
\textbf{prononciation à la française} --- dites-la à voix haute en lisant le tableau.
Créez vos cartes Anki au fur et à mesure, 15 mots par jour maximum.

\subsection{Thème 1 --- Se présenter, identité, état civil}
\begin{tabvocab}
\mot{first name}{feust né-im}{prénom --- \en{What's your first name?}}
\mot{last name / surname}{laast né-im}{nom de famille}
\mot{full name}{foul né-im}{nom complet}
\mot{date of birth}{dé-it euv beuss}{date de naissance}
\mot{place of birth}{pléiss euv beuss}{lieu de naissance}
\mot{age}{é-idj}{âge --- \en{How old are you?} / \en{I'm 30.}}
\mot{nationality}{na-cheu-na-LI-ti}{nationalité --- \en{Senegalese, French, Nigerian}}
\mot{address}{e-DRESS}{adresse (accent sur la 2\textsuperscript{e} syllabe !)}
\mot{phone number}{foun NEUM-beu}{numéro de téléphone}
\mot{email}{ii-mé-il}{e-mail --- \en{email address}}
\mot{marital status}{MA-ri-tl STÉ-i-teuss}{situation matrimoniale}
\mot{single}{si-ngl}{célibataire}
\mot{married}{MA-rid}{marié(e)}
\mot{divorced}{di-VOOSST}{divorcé(e)}
\mot{widow / widower}{ouii-dou}{veuve / veuf}
\mot{husband}{HOUZ-beund}{mari}
\mot{wife}{ouaïf}{femme (épouse) --- pluriel : \en{wives}}
\mot{job}{djob}{emploi, métier --- \en{What's your job?}}
\mot{occupation}{o-kieu-PÉ-i-cheun}{profession (formel)}
\mot{student}{STiou-d'nt}{étudiant, élève}
\mot{pupil}{piou-pl}{élève (école primaire, BrE)}
\mot{ID card}{aï-di kaad}{carte d'identité}
\mot{passport}{PAASS-poort}{passeport}
\mot{signature}{SIG-neu-tcheu}{signature}
\mot{to spell}{spèl}{épeler --- \en{How do you spell your name?}}
\end{tabvocab}

\subsection{Thème 2 --- La famille et les relations}
\begin{tabvocab}
\mot{parents}{PÉ-e-reunts}{parents --- \en{my mother and father}}
\mot{mother / mom / mum}{MO-dheu}{mère, maman}
\mot{father / dad}{FAA-dheu}{père, papa}
\mot{son}{seun}{fils}
\mot{daughter}{DOO-teu}{fille (enfant) --- le \en{gh} est muet}
\mot{child / children}{tchaïld / TCHIL-dren}{enfant / enfants}
\mot{brother}{BRA-dheu}{frère}
\mot{sister}{SI-steu}{s\oe ur}
\mot{uncle}{EUN-kl}{oncle}
\mot{aunt}{aant}{tante}
\mot{cousin}{KEU-z'n}{cousin, cousine}
\mot{nephew}{NÉ-fiou}{neveu}
\mot{niece}{niiss}{nièce}
\mot{grandmother}{GRAND-mo-dheu}{grand-mère}
\mot{grandfather}{GRAND-faa-dheu}{grand-père}
\mot{grandson / granddaughter}{GRAND-seun}{petit-fils / petite-fille}
\mot{in-laws}{in-looz}{la belle-famille --- \en{mother-in-law, brother-in-law}}
\mot{stepfather / stepmother}{STÈP-faa-dheu}{beau-père / belle-mère (remariage)}
\mot{half-brother}{haaf-BRA-dheu}{demi-frère}
\mot{twin}{touin}{jumeau --- \en{twin brother}}
\mot{relative}{RÉ-leu-tiv}{membre de la famille, parent}
\mot{friend}{frènd}{ami --- \en{best friend, close friend}}
\mot{neighbour}{NÉ-i-beu}{voisin}
\mot{colleague}{KO-liig}{collègue}
\mot{boyfriend / girlfriend}{BOÏ-frend}{petit ami / petite amie}
\mot{fiancé(e)}{fi-ON-sé-i}{fiancé(e)}
\mot{to get married}{guèt MA-rid}{se marier}
\mot{to give birth}{guiv beuss}{accoucher, donner naissance}
\mot{to bring up / raise}{bri-ng UP}{élever (des enfants)}
\mot{to take care of}{té-ik KÉ-eu euv}{s'occuper de}
\end{tabvocab}

\subsection{Thème 3 --- Le corps, la santé, le médecin}
\begin{tabvocab}
\mot{head}{hèd}{tête --- \en{I have a headache}}
\mot{hair}{hè-eu}{cheveux (indénombrable) --- \en{long hair}}
\mot{face}{fé-iss}{visage}
\mot{eye / eyes}{aï}{\oe il / yeux --- \en{brown eyes}}
\mot{ear}{i-eu}{oreille}
\mot{nose}{nôouz}{nez}
\mot{mouth}{maouth}{bouche --- pluriel \en{mouths} \pron{maou-dhz}}
\mot{tooth / teeth}{touuth / tiith}{dent / dents}
\mot{tongue}{teung}{langue (organe)}
\mot{lip}{lip}{lèvre}
\mot{neck}{nèk}{cou}
\mot{shoulder}{CHOUL-deu}{épaule}
\mot{arm}{aam}{bras}
\mot{hand}{hand}{main}
\mot{finger}{FI-ngheu}{doigt}
\mot{chest}{tchèst}{poitrine}
\mot{back}{bak}{dos}
\mot{stomach}{STO-meuk}{ventre, estomac}
\mot{leg}{lèg}{jambe}
\mot{knee}{nii}{genou --- le \en{k} est muet}
\mot{foot / feet}{fout / fiit}{pied / pieds}
\mot{skin}{skine}{peau}
\mot{blood}{bleud}{sang}
\mot{bone}{bôoun}{os}
\mot{heart}{haat}{c\oe ur (organe) --- \en{heart attack}}
\mot{ill / sick}{il / sik}{malade}
\mot{pain}{pé-in}{douleur --- \en{a pain in my back}}
\mot{headache}{HÈ-dé-ik}{mal de tête}
\mot{stomachache}{STO-meuk-é-ik}{mal de ventre}
\mot{fever / temperature}{FII-veu}{fièvre --- \en{I have a fever}}
\mot{cough}{kof}{toux --- \en{to cough} \pron{kof}}
\mot{cold}{kôould}{rhume --- \en{to catch a cold} ; aussi : froid}
\mot{flu}{flou}{grippe}
\mot{sore throat}{soo THRO-OUT}{mal de gorge}
\mot{runny nose}{RA-ni nôouz}{nez qui coule}
\mot{dizzy}{DI-zi}{qui tourne la tête --- \en{I feel dizzy}}
\mot{tired}{TAÏ-eud}{fatigué}
\mot{injury / wound}{IN-djeu-ri}{blessure}
\mot{to bleed}{bliid}{saigner}
\mot{medicine}{MÉ-deu-sine}{médicament, médecine}
\mot{pill / tablet}{pil}{comprimé}
\mot{prescription}{pri-SKRIP-cheun}{ordonnance}
\mot{pharmacy / chemist}{FAA-meu-si}{pharmacie}
\mot{doctor / GP}{DOOK-teu}{médecin --- \en{general practitioner}}
\mot{nurse}{neuss}{infirmier, infirmière}
\mot{hospital}{HOSS-pi-tl}{hôpital}
\mot{appointment}{e-POÏNT-m'nt}{rendez-vous --- \en{to make an appointment}}
\mot{surgery}{SEU-djeu-ri}{cabinet médical (BrE) ; chirurgie}
\mot{treatment}{TRIIT-m'nt}{traitement}
\mot{vaccine}{VA-siinn}{vaccin}
\mot{health insurance}{hèlth in-Chou-reuns}{assurance maladie}
\mot{to recover}{ri-KEU-veu}{guérir, se rétablir}
\end{tabvocab}

\begin{astuce}[Les dix phrases médicales qui sauvent]
\en{I don't feel well.} (Je ne me sens pas bien.) \quad
\en{I have a headache / a fever / a sore throat.} \quad
\en{It hurts here.} (Ça fait mal ici.) \quad
\en{I'm allergic to\ldots} \quad
\en{I've been feeling ill for three days.} (Je suis malade depuis trois jours.) \quad
\en{I need to see a doctor.} \quad
\en{Where is the nearest pharmacy?} \quad
\en{Is it serious?} \quad
\en{What should I do?}
\end{astuce}

\subsection{Thème 4 --- La maison et les objets du quotidien}
\begin{tabvocab}
\mot{house / home}{haouss / hâam}{maison (bâtiment) / chez soi}
\mot{flat / apartment}{flatt / e-PAAUT-m'nt}{appartement (BrE / AmE)}
\mot{room}{rouum}{pièce}
\mot{living room}{LI-ving rouum}{salon}
\mot{bedroom}{BÈD-roum}{chambre}
\mot{kitchen}{KI-tchin}{cuisine}
\mot{bathroom}{BAAth-roum}{salle de bains}
\mot{toilet}{TO-i-l't}{toilettes --- \en{Where is the toilet?}}
\mot{hall / entrance}{hooll}{entrée, couloir}
\mot{balcony}{BA-l-keu-ni}{balcon}
\mot{yard / garden}{iaad / GAA-d'n}{cour / jardin}
\mot{roof}{rouuf}{toit}
\mot{wall}{ouool}{mur}
\mot{floor}{floo}{sol, étage --- \en{on the ground floor}}
\mot{ceiling}{SII-ling}{plafond}
\mot{door}{doo}{porte}
\mot{window}{OUIN-dou}{fenêtre}
\mot{stairs}{stè-euz}{escalier --- \en{go up the stairs}}
\mot{lift / elevator}{lift / É-leu-vé-i-teu}{ascenseur (BrE / AmE)}
\mot{key}{kii}{clé}
\mot{lock}{lok}{serrure ; \en{to lock} verrouiller}
\mot{furniture}{FEU-ni-tcheu}{meubles (indénombrable)}
\mot{table}{TÉ-i-bl}{table}
\mot{chair}{tchè-eu}{chaise}
\mot{armchair}{AArm-tchè-eu}{fauteuil}
\mot{sofa / couch}{SÔ-fa}{canapé}
\mot{bed}{bèd}{lit --- \en{to make the bed}}
\mot{cupboard}{KEU-beud}{placard --- le \en{p} est muet}
\mot{shelf / shelves}{chèlf}{étagère}
\mot{desk}{dèsk}{bureau (meuble)}
\mot{mirror}{MI-reu}{miroir}
\mot{lamp}{lamp}{lampe}
\mot{light / lamp}{laït}{lumière / luminaire}
\mot{fridge}{fridj}{réfrigérateur --- \en{refrigerator}}
\mot{freezer}{FRII-zeu}{congélateur}
\mot{oven}{EUV'n}{four}
\mot{stove / cooker}{stôouv}{cuisinière}
\mot{kettle}{KÈ-tl}{bouilloire}
\mot{sink}{sink}{évier}
\mot{tap / faucet}{tap}{robinet (BrE / AmE)}
\mot{washing machine}{OUO-ching me-chiin}{machine à laver}
\mot{iron}{AÏ-eun}{fer à repasser --- aussi : fer (métal)}
\mot{broom}{brouum}{balai}
\mot{bucket}{BEU-k't}{seau}
\mot{towel}{TA-ouel}{serviette}
\mot{soap}{sôoup}{savon}
\mot{blanket}{BLANG-k't}{couverture}
\mot{pillow}{PI-lou}{oreiller}
\mot{sheet}{chiit}{drap ; feuille}
\mot{electricity}{i-lèk-TRI-si-ti}{électricité}
\mot{water}{OUO-teu}{eau}
\mot{power cut}{PA-ou-eu keut}{coupure de courant}
\end{tabvocab}

\subsection{Thème 5 --- La nourriture et les boissons}
\begin{tabvocab}
\mot{food}{fouud}{nourriture}
\mot{meal}{miil}{repas --- \en{three meals a day}}
\mot{breakfast}{BRÈK-feust}{petit-déjeuner --- \en{to have breakfast}}
\mot{lunch}{leuntch}{déjeuner}
\mot{dinner}{DI-neu}{dîner}
\mot{snack}{snak}{en-cas, casse-croûte}
\mot{bread}{brèd}{pain --- \en{a loaf of bread}}
\mot{rice}{raïss}{riz}
\mot{meat}{miit}{viande}
\mot{beef}{biif}{b\oe uf}
\mot{chicken}{TCHI-k'n}{poulet}
\mot{fish}{fich}{poisson --- pluriel \en{fish}}
\mot{egg}{èg}{\oe uf}
\mot{vegetable}{VEDJ-teu-bl}{légume}
\mot{onion}{EUN-ieun}{oignon}
\mot{tomato}{teu-MÉ-i-tou}{tomate --- pluriel \en{tomatoes}}
\mot{potato}{peu-TÉ-i-tou}{pomme de terre}
\mot{carrot}{KA-rot}{carotte}
\mot{cabbage}{KA-bidj}{chou}
\mot{salad}{SA-lad}{salade}
\mot{fruit}{frouut}{fruit}
\mot{banana}{be-NA-neu}{banane}
\mot{orange}{O-rindj}{orange}
\mot{mango}{MANG-gou}{mangue}
\mot{lemon}{LÉ-meun}{citron}
\mot{apple}{A-pl}{pomme}
\mot{sugar}{CHOU-geu}{sucre}
\mot{salt}{soolt}{sel}
\mot{pepper}{PÉ-peu}{poivre}
\mot{oil}{oïl}{huile}
\mot{butter}{BEU-teu}{beurre}
\mot{cheese}{tchiiz}{fromage}
\mot{milk}{milk}{lait}
\mot{yoghurt}{IO-gheuut}{yaourt}
\mot{water}{OUO-teu}{eau --- \en{still / sparkling}}
\mot{tea / coffee}{tii / KO-fi}{thé / café}
\mot{juice}{djouuss}{jus}
\mot{soda / soft drink}{SÔ-da}{boisson gazeuse}
\mot{sauce}{sooss}{sauce}
\mot{soup}{souup}{soupe}
\mot{dish}{dich}{plat --- \en{a traditional dish}}
\mot{recipe}{RÉ-si-pi}{recette}
\mot{ingredients}{in-GRI-di-eunts}{ingrédients}
\mot{to cook}{kouk}{cuisiner}
\mot{to boil}{boïl}{bouillir}
\mot{to fry}{fraï}{frire}
\mot{to grill}{gril}{griller}
\mot{to bake}{bé-ik}{cuire au four}
\mot{to roast}{rouust}{rôtir}
\mot{to cut}{keut}{couper}
\mot{to mix}{miks}{mélanger}
\mot{to peel}{piil}{éplucher}
\mot{to taste}{té-ist}{goûter ; avoir le goût de}
\mot{delicious}{di-LI-cheus}{délicieux}
\mot{tasty}{TÉ-i-sti}{savoureux}
\mot{spicy}{SPA-i-si}{épicé}
\mot{sweet}{souit}{sucré}
\mot{sour}{SA-ou-eu}{acide, aigre}
\mot{bitter}{BI-teu}{amer}
\mot{fresh}{frèch}{frais}
\mot{raw}{roo}{cru}
\mot{ripe}{raïp}{mûr}
\mot{hungry}{HEU-ngri}{avoir faim --- \en{I'm hungry}}
\mot{thirsty}{THEU-sti}{avoir soif}
\mot{full}{foul}{plein, repu --- \en{I'm full}}
\mot{leftovers}{LÈFT-ou-veuz}{restes (du repas)}
\end{tabvocab}

\subsection{Thème 6 --- Au restaurant et au marché}
\begin{tabvocab}
\mot{restaurant}{RÈSS-teu-ront}{restaurant}
\mot{waiter / waitress}{OUÉ-i-teu}{serveur / serveuse}
\mot{menu}{MÉ-niou}{carte, menu}
\mot{bill / check}{bil / tchèk}{addition (BrE / AmE)}
\mot{tip}{tip}{pourboire}
\mot{reservation / booking}{rè-zeu-VÉ-i-cheun}{réservation --- \en{to book a table}}
\mot{order}{OO-deu}{commande ; \en{to order} commander}
\mot{dish of the day}{dich euv dhe DÉ-i}{plat du jour}
\mot{starter / appetizer}{STAA-teu}{entrée}
\mot{main course}{mé-in kooss}{plat principal}
\mot{dessert}{di-ZEUT}{dessert}
\mot{drink}{drink}{boisson}
\mot{glass / bottle}{glaass / BO-tl}{verre / bouteille}
\mot{plate}{plé-it}{assiette}
\mot{bowl}{bôoul}{bol}
\mot{fork / knife / spoon}{fook / naïf / spouun}{fourchette / couteau / cuillère}
\mot{napkin}{NA-pkine}{serviette de table}
\mot{market}{MAA-k't}{marché}
\mot{shop / store}{chop / stoo}{magasin}
\mot{supermarket}{SOU-peu-maa-k't}{supermarché}
\mot{customer}{KEUSS-teu-meu}{client}
\mot{seller}{SÉ-leu}{vendeur}
\mot{price}{praïss}{prix}
\mot{cheap / expensive}{tchiip / ik-SPÈN-siv}{bon marché / cher}
\mot{discount}{DIS-kaount}{remise, réduction}
\mot{bargain}{BAA-g'n}{bonne affaire}
\mot{receipt}{ri-SIIT}{reçu, ticket --- le \en{p} est muet}
\mot{cash}{kach}{espèces --- \en{to pay in cash}}
\mot{change}{tché-indj}{monnaie ; changement}
\mot{shopping}{CHO-ping}{courses --- \en{to do the shopping}}
\mot{basket / trolley}{BAAS-k't}{panier / chariot}
\mot{queue}{kiou}{file d'attente --- \en{to queue up}}
\mot{size}{saïz}{taille, pointure}
\mot{weigh}{oué-i}{peser}
\mot{kilo / gram}{KII-lou}{kilo / gramme}
\end{tabvocab}

\begin{astuce}[Les 12 phrases du restaurant]
\en{A table for two, please.} \quad \en{Could I see the menu?} \quad
\en{What do you recommend?} \quad \en{I'd like the grilled fish.} \quad
\en{I'll have the same.} \quad \en{Could I have some water, please?} \quad
\en{It's delicious, thank you.} \quad \en{I'm vegetarian / I don't eat pork.} \quad
\en{Could we have the bill, please?} \quad \en{Is service included?} \quad
\en{Keep the change.} \quad \en{Everything was lovely, thanks.}
\end{astuce}

\subsection{Thème 7 --- Les vêtements et les couleurs}
\begin{tabvocab}
\mot{clothes}{klôoudhz}{vêtements (toujours pluriel)}
\mot{shirt}{cheut}{chemise}
\mot{T-shirt}{TII-cheut}{t-shirt}
\mot{trousers / pants}{TRA-ou-zeuz}{pantalon (BrE / AmE)}
\mot{jeans}{djiinz}{jean}
\mot{shorts}{choo-ts}{short}
\mot{skirt}{skeut}{jupe}
\mot{dress}{drèss}{robe}
\mot{suit}{sout}{costume --- aussi : \en{to suit} aller à (qqn)}
\mot{jacket}{DJA-k't}{veste}
\mot{coat}{kôout}{manteau}
\mot{sweater / jumper}{SOÉ-teu}{pull}
\mot{hoodie}{HOU-di}{sweat à capuche}
\mot{shoes}{chouuz}{chaussures}
\mot{trainers / sneakers}{TRÉ-i-neuz}{baskets}
\mot{sandals}{SAN-deulz}{sandales}
\mot{boots}{bouuts}{bottes}
\mot{slippers}{SLI-peuz}{pantoufles}
\mot{socks}{soks}{chaussettes}
\mot{hat / cap}{hatt / kap}{chapeau / casquette}
\mot{scarf}{skaaf}{écharpe --- pluriel \en{scarves}}
\mot{gloves}{gleuvz}{gants}
\mot{belt}{bèlt}{ceinture}
\mot{bag}{bag}{sac}
\mot{purse / wallet}{peuss / OUO-l't}{porte-monnaie / portefeuille}
\mot{glasses}{glaa-siz}{lunettes}
\mot{watch}{ouotch}{montre}
\mot{jewellery / jewelry}{DJO-ou-el-ri}{bijoux}
\mot{ring}{ring}{bague ; anneau}
\mot{umbrella}{em-BRÉ-la}{parapluie}
\mot{to wear}{oué-eu}{porter (vêtement) --- \en{She wears a dress}}
\mot{to put on}{pout ON}{mettre (action)}
\mot{to take off}{té-ik OF}{enlever}
\mot{to try on}{traï ON}{essayer}
\mot{to fit}{fit}{aller à la bonne taille}
\mot{size}{saïz}{taille}
\mot{red, blue, green, yellow}{rèd, blou, griin, IÉ-lou}{rouge, bleu, vert, jaune}
\mot{black, white, grey / gray}{blak, ouaït, gré-i}{noir, blanc, gris}
\mot{brown, orange, pink, purple}{braoun, O-rindj, pink, PEU-pl}{marron, orange, rose, violet}
\mot{light blue / dark blue}{laït blou}{bleu clair / bleu foncé}
\mot{colourful}{KEU-leu-ful}{coloré}
\end{tabvocab}

\subsection{Thème 8 --- La ville, les transports, se déplacer}
\begin{tabvocab}
\mot{city / town}{SI-ti / taoun}{grande ville / ville}
\mot{village}{VI-lidj}{village}
\mot{street}{striit}{rue --- \en{main street}}
\mot{road}{rôoud}{route}
\mot{avenue / boulevard}{A-veu-niou}{avenue}
\mot{square}{skoué-eu}{place}
\mot{corner}{KOO-neu}{coin --- \en{at the corner}}
\mot{crossroads}{KROSS-rôoudz}{carrefour}
\mot{traffic lights}{TRA-fik laïts}{feux de circulation}
\mot{roundabout}{RAUND-eu-baout}{rond-point}
\mot{bridge}{bridj}{pont}
\mot{market}{MAA-k't}{marché}
\mot{mosque / church}{mosk / tcheutch}{mosquée / église}
\mot{school}{skouul}{école}
\mot{hospital}{HOSS-pi-tl}{hôpital}
\mot{bank}{bangk}{banque}
\mot{post office}{pôoust O-fiss}{bureau de poste}
\mot{police station}{peu-LIISS STÉ-i-cheun}{commissariat}
\mot{petrol station / gas station}{PÉ-tr'l}{station-service}
\mot{bus stop}{beuss stop}{arrêt de bus}
\mot{bus / coach}{beuss}{bus / autocar}
\mot{car}{kaa}{voiture}
\mot{taxi}{TA-k-si}{taxi}
\mot{train}{tré-in}{train}
\mot{plane / airplane}{plé-in}{avion}
\mot{boat / ship}{bôout / chip}{bateau / navire}
\mot{bicycle / motorbike}{BA-i-sik'l}{vélo / moto}
\mot{ticket}{TI-k't}{billet, ticket}
\mot{fare}{fé-eu}{tarif, prix du trajet}
\mot{platform}{PLAT-foom}{quai, plateforme}
\mot{station / airport}{STÉ-i-cheun}{gare / aéroport}
\mot{traffic}{TRA-fik}{circulation}
\mot{traffic jam}{TRA-fik djam}{embouteillage}
\mot{driver}{DRAÏ-veu}{conducteur}
\mot{passenger}{PA-ss'n-djeu}{passager}
\mot{to drive}{draïv}{conduire}
\mot{to ride}{raïd}{rouler, monter (à vélo, à cheval)}
\mot{to walk}{ouook}{marcher --- \en{on foot}}
\mot{to take the bus}{té-ik dhe beuss}{prendre le bus}
\mot{to get on / off}{guèt ON / OF}{monter / descendre (bus, train)}
\mot{to get in / out of}{guèt IN / AOUT euv}{monter / sortir (voiture)}
\mot{to turn left / right}{teun lèft / raït}{tourner à gauche / à droite}
\mot{to go straight on}{gou stré-it ON}{aller tout droit}
\mot{to cross}{kross}{traverser}
\mot{opposite}{O-ppo-zit}{en face de}
\mot{near / far from}{ni-eu / faa}{près / loin de}
\mot{next to}{nèkst tou}{à côté de}
\mot{between}{bi-TOUINN}{entre}
\mot{around / about}{e-RAOUND}{autour de, environ}
\mot{direction}{di-RÈK-cheun}{direction}
\mot{map}{map}{carte, plan}
\mot{distance}{DI-ss'teuns}{distance --- \en{It's two kilometres away.}}
\end{tabvocab}

\subsection{Thème 9 --- Le temps, les routines, le calendrier}
\begin{tabvocab}
\mot{time}{taïm}{temps, heure, fois}
\mot{clock / watch}{klok / ouotch}{horloge / montre}
\mot{hour / minute / second}{A-ou-eu / MI-nit / SÉ-keund}{heure / minute / seconde}
\mot{day / week / month / year}{dé-i / ouii-k / meunth / i-eu}{jour / semaine / mois / année}
\mot{today / tomorrow / yesterday}{teu-DÉ-i / te-MO-rou / IÈS-teu-dé-i}{aujourd'hui / demain / hier}
\mot{morning / afternoon / evening}{MOO-ning / aaf-teu-NOOUN / II-v'-ning}{matin / après-midi / soir}
\mot{night / midnight / noon}{naït / MID-naït / nouun}{nuit / minuit / midi}
\mot{Monday \ldots Sunday}{MEUN-di \ldots SEUN-di}{lundi \ldots dimanche}
\mot{January \ldots December}{DJA-niou-é-ri \ldots di-SSÈM-beu}{janvier \ldots décembre}
\mot{season}{SII-z'n}{saison --- \en{dry season, rainy season}}
\mot{weekend}{ouii-k'ènd}{week-end}
\mot{holiday / vacation}{HO-li-dé-i}{vacances, jour férié}
\mot{daily routine}{DÉ-i-li rou-TIIN}{routine quotidienne}
\mot{to wake up / get up}{oué-ik EP / guèt EP}{se réveiller / se lever}
\mot{to wash / shower}{ouoch / CHA-ou-eu}{se laver / se doucher}
\mot{to get dressed}{guèt drèst}{s'habiller}
\mot{to pray}{pré-i}{prier}
\mot{to commute}{keu-MIOUT}{faire le trajet domicile-travail}
\mot{to work / study}{oueu-k / STOU-di}{travailler / étudier}
\mot{to rest / relax}{rèst / ri-LAKS}{se reposer / se détendre}
\mot{to go to bed}{gou tou bèd}{aller se coucher}
\mot{to sleep}{sliip}{dormir --- \en{slept, slept}}
\mot{early / late}{EU-li / lé-it}{tôt / tard ; en retard}
\mot{on time / in time}{on taïm}{à l'heure / à temps}
\mot{busy / free}{BI-zi / frii}{occupé / libre}
\mot{always / never}{AAL-oué-iz / NÉ-veu}{toujours / jamais}
\end{tabvocab}

\subsection{Thème 10 --- Le travail et les métiers}
\begin{tabvocab}
\mot{work / job}{oueu-k / djob}{travail (activité) / emploi (poste)}
\mot{company / firm}{KEUM-peu-ni}{entreprise}
\mot{office}{O-fiss}{bureau (lieu)}
\mot{boss / manager}{boss / MA-n'e-djeu}{patron / directeur}
\mot{employee / employer}{em-PLOÏ-ii / em-PLOÏ-eu}{employé / employeur}
\mot{staff / team}{staaf / tiim}{personnel / équipe}
\mot{colleague}{KO-liig}{collègue}
\mot{customer / client}{KEUSS-teu-meu}{client}
\mot{salary / wage}{SA-leu-ri / oué-idj}{salaire (mensuel) / salaire (horaire)}
\mot{income}{IN-keum}{revenu}
\mot{contract}{KON-trakt}{contrat}
\mot{interview}{IN-teu-viou}{entretien d'embauche}
\mot{CV / resume}{si-vi / RÈ-ziou-mé-i}{curriculum vitæ}
\mot{experience}{ik-SPÍ-e-ri-euns}{expérience}
\mot{skill}{skil}{compétence}
\mot{training}{TRÉ-i-ning}{formation}
\mot{meeting}{MII-ting}{réunion}
\mot{deadline}{DÈD-laïn}{date limite}
\mot{task / project}{taask / PRO-djèkt}{tâche / projet}
\mot{shift}{chift}{poste, horaire --- \en{night shift}}
\mot{overtime}{Ô-veu-taïm}{heures supplémentaires}
\mot{full-time / part-time}{foul-taïm}{à temps plein / partiel}
\mot{unemployed}{eun-im-PLOÏD}{au chômage}
\mot{to hire / employ}{haï-eu}{embaucher}
\mot{to fire / sack}{faï-eu}{licencier}
\mot{to resign}{ri-ZAÏN}{démissionner}
\mot{to retire}{ri-TAÏ-eu}{prendre sa retraite}
\mot{to earn}{eun}{gagner (de l'argent)}
\mot{to apply for}{e-PLAÏ fo}{postuler à}
\mot{teacher, doctor, nurse}{TII-tcheu}{enseignant, médecin, infirmier}
\mot{engineer, technician}{en-dji-NI-eu}{ingénieur, technicien}
\mot{driver, mechanic}{DRAÏ-veu, me-KA-nik}{chauffeur, mécanicien}
\mot{farmer, fisherman}{FAA-meu}{agriculteur, pêcheur}
\mot{trader, shopkeeper}{TRÉ-i-deu}{commerçant}
\mot{tailor, carpenter, mason}{TÉ-i-lou}{tailleur, menuisier, maçon}
\mot{civil servant}{SI-v'l SEU-veunt}{fonctionnaire}
\mot{lawyer, judge}{LOÏ-eu}{avocat, juge}
\mot{police officer}{peu-LIISS O-fi-seu}{policier}
\mot{soldier}{SÔUL-djeu}{soldat}
\mot{accountant}{e-KAUN-teunt}{comptable}
\mot{secretary}{SÉ-creu-teu-ri}{secrétaire}
\mot{chef / cook}{chèf / kouk}{chef cuisinier / cuisinier}
\mot{hairdresser}{HÉ-eu-drè-seu}{coiffeur}
\mot{journalist}{DJEU-neu-list}{journaliste}
\end{tabvocab}

\subsection{Thème 11 --- L'école, les études, l'apprentissage}
\begin{tabvocab}
\mot{school}{skouul}{école}
\mot{primary / secondary school}{PRAÏ-me-ri}{école primaire / secondaire}
\mot{university}{iou-ni-VEU-si-ti}{université}
\mot{class / classroom}{klaass}{classe, cours}
\mot{lesson / course}{LÈ-ss'n / kooss}{leçon / cours (AmE : cours ; BrE : \en{course} = formation)}
\mot{subject}{SEUB-djèkt}{matière}
\mot{teacher / lecturer}{TII-tcheu}{enseignant / professeur d'université}
\mot{pupil / student}{piou-pl / STiou-d'nt}{élève / étudiant}
\mot{exam / test}{ig-ZAM / tèst}{examen / contrôle}
\mot{result / grade / mark}{ri-ZEULT}{résultat / note}
\mot{to pass / fail an exam}{paass / fé-il}{réussir / échouer à un examen}
\mot{homework}{HÔUM-oueu-k}{devoirs (à la maison)}
\mot{exercise}{ÉK-seu-saïz}{exercice}
\mot{question / answer}{KWÈS-tchoun / AA-n'seu}{question / réponse}
\mot{dictionary}{DIK-cheu-neu-ri}{dictionnaire}
\mot{notebook / textbook}{NÔUT-bouk}{cahier / manuel}
\mot{pen, pencil, eraser}{pèn, PÈN-sl, i-RÉ-i-zeu}{stylo, crayon, gomme}
\mot{board}{bood}{tableau --- \en{whiteboard, blackboard}}
\mot{to learn}{leun}{apprendre (acquérir)}
\mot{to study}{STOU-di}{étudier (travailler)}
\mot{to teach}{tiitch}{enseigner --- \en{taught, taught}}
\mot{to explain}{ik-SPLÉ-IN}{expliquer}
\mot{to understand}{en-deu-STAND}{comprendre --- \en{understood}}
\mot{to remember / forget}{ri-MÈM-beu / feu-GÈT}{se souvenir / oublier}
\mot{to practise}{PRAK-tiss}{s'entraîner}
\mot{to improve}{im-PROUUV}{améliorer}
\mot{to graduate}{GRA-djou-é-it}{obtenir son diplôme}
\mot{degree / diploma}{di-GRIi}{diplôme universitaire / diplôme}
\mot{scholarship}{SKO-leu-chip}{bourse d'études}
\mot{knowledge}{NO-lidj}{connaissances, savoir}
\mot{skill}{skil}{compétence}
\mot{progress}{PRO-grès}{progrès --- \en{to make progress}}
\end{tabvocab}

\subsection{Thème 12 --- L'argent, la banque, les achats en ligne}
\begin{tabvocab}
\mot{money}{MEU-ni}{argent (indénombrable)}
\mot{cash}{kach}{espèces}
\mot{coin / banknote}{koïn / BANGK-nôut}{pièce / billet}
\mot{currency}{KEU-ren-si}{monnaie --- \en{the CFA franc, the dollar, the euro}}
\mot{exchange rate}{iks-TCHÉ-indj ré-it}{taux de change}
\mot{bank}{bangk}{banque}
\mot{bank account}{bang e-KAOUNT}{compte bancaire}
\mot{credit / debit card}{KRÉ-dit kaad}{carte de crédit / débit}
\mot{ATM / cash machine}{é-i-tii-èm}{distributeur automatique}
\mot{transfer}{TRANS-feu}{virement}
\mot{loan}{lôoun}{prêt}
\mot{debt}{dèt}{dette --- le \en{b} est muet}
\mot{insurance}{in-CHOU-reuns}{assurance}
\mot{tax / duty}{taks / DIOUTI}{impôt / taxe}
\mot{receipt}{ri-SIIT}{reçu}
\mot{invoice}{IN-voïs}{facture}
\mot{budget}{BEU-djit}{budget}
\mot{savings}{SÉ-i-vingz}{économies}
\mot{to pay}{pé-i}{payer --- \en{paid, paid}}
\mot{to spend}{spènd}{dépenser --- \en{spent, spent}}
\mot{to save}{sé-iv}{économiser ; sauver}
\mot{to cost}{kost}{coûter --- \en{cost, cost}}
\mot{to borrow / lend}{BO-rou / lènd}{emprunter / prêter}
\mot{to owe}{ôou}{devoir (de l'argent)}
\mot{to refund}{ri-FEUND}{rembourser}
\mot{to afford}{e-FOOD}{avoir les moyens de}
\mot{online shopping}{ON-laïn CHO-ping}{achats en ligne}
\mot{to order / deliver}{OO-deu / di-LI-veu}{commander / livrer}
\mot{delivery fee}{di-LI-ve-ri fii}{frais de livraison}
\end{tabvocab}

\subsection{Thème 13 --- Le voyage, l'aéroport, l'hôtel}
\begin{tabvocab}
\mot{trip / journey}{trip / DJEU-ni}{voyage (court) / trajet}
\mot{travel}{TRA-v'l}{voyager (verbe)}
\mot{flight}{flaït}{vol --- \en{flight number}}
\mot{airport}{É-eu-poort}{aéroport}
\mot{check-in}{TCHÈK-in}{enregistrement}
\mot{boarding pass}{BOO-ding paass}{carte d'embarquement}
\mot{luggage / baggage}{LEU-gidj}{bagages}
\mot{suitcase}{SOUT-ké-iss}{valise}
\mot{hand luggage}{hand LEU-gidj}{bagage à main}
\mot{passport / visa}{PAASS-poort / VII-za}{passeport / visa}
\mot{customs}{KEUSS-teumz}{douane}
\mot{security check}{si-KIO-ou-ri-ti}{contrôle de sécurité}
\mot{gate}{gé-it}{porte d'embarquement}
\mot{departure / arrival}{di-PAA-tcheu / e-RAÏ-veul}{départ / arrivée}
\mot{delay}{di-LÉ-i}{retard}
\mot{to book}{bouk}{réserver}
\mot{hotel}{hÔ-TÈL}{hôtel}
\mot{reception / front desk}{ri-SÈP-cheun}{réception}
\mot{single / double room}{SI-ngl / DOU-bl roum}{chambre simple / double}
\mot{key / key card}{kii}{clé / carte-clé}
\mot{floor}{floo}{étage}
\mot{air conditioning}{É-eu keun-di-CHEU-ning}{climatisation}
\mot{breakfast included}{BRÈK-feust in-KLOU-did}{petit-déjeuner compris}
\mot{tourist / guide}{TO-ou-rist / gaïd}{touriste / guide}
\mot{sightseeing}{SAÏT-sii-ing}{visites touristiques}
\mot{beach}{biitch}{plage}
\mot{embassy}{ÈM-beu-si}{ambassade}
\end{tabvocab}

\subsection{Thème 14 --- La nature, la météo, les animaux}
\begin{tabvocab}
\mot{weather}{OUÉ-dheu}{temps (météo) --- \en{What's the weather like?}}
\mot{sunny / cloudy}{SEU-ni / KLA-ou-di}{ensoleillé / nuageux}
\mot{hot / warm / cool / cold}{hot / ouoom / kouul / kôould}{chaud / tiède / frais / froid}
\mot{rain / rainy}{ré-in}{pluie --- \en{It's raining}}
\mot{storm / thunder}{stoom / THEU-n'deu}{tempête / tonnerre}
\mot{wind / windy}{ouind}{vent --- \en{It's windy}}
\mot{humidity}{hiou-MI-di-ti}{humidité}
\mot{heat / drought}{hiit / draout}{chaleur / sécheresse}
\mot{flood}{fleud}{inondation}
\mot{sky / sun / moon / star}{skaï / seun / mouun / staa}{ciel / soleil / lune / étoile}
\mot{sea / ocean / river / lake}{sii / Ô-cheun / RI-veu / lé-ik}{mer / océan / fleuve / lac}
\mot{mountain / hill}{MAUN-t'n / hil}{montagne / colline}
\mot{forest / tree / plant}{FO-rist / trii / plaant}{forêt / arbre / plante}
\mot{sand / stone / soil}{sand / stôoun / soïl}{sable / pierre / sol}
\mot{animal}{A-ni-ml}{animal --- pluriel \en{animals}}
\mot{dog / cat / horse / cow / sheep / goat}{dog / katt / hooss / kaou / chiip / gôut}{chien / chat / cheval / vache / mouton / chèvre}
\mot{chicken / duck / bird}{TCHI-k'n / deuk / beud}{poule / canard / oiseau}
\mot{lion / elephant / monkey / snake}{LAÏ-eun}{lion / éléphant / singe / serpent}
\mot{fish / mosquito / fly / spider}{fich / meuss-KII-tou / flaï / SPAÏ-deu}{poisson / moustique / mouche / araignée}
\mot{tail / wing / paw}{té-il / ouing / poo}{queue / aile / patte}
\end{tabvocab}

\subsection{Thème 15 --- Loisirs, sport, musique, cinéma}
\begin{tabvocab}
\mot{hobby}{HO-bi}{passe-temps --- pluriel \en{hobbies}}
\mot{free time}{frii taïm}{temps libre}
\mot{sport / match / game}{spoort / mattch / gé-im}{sport / match / jeu}
\mot{football / basketball}{FOOT-booll}{football / basketball}
\mot{team / player / coach}{tiim / PLÉ-i-eu / kôoutch}{équipe / joueur / entraîneur}
\mot{to win / lose / draw}{ouine / louz / droo}{gagner / perdre / faire match nul}
\mot{goal / score}{gôoul / skoo}{but / score}
\mot{championship / tournament}{TCHAM-pieun-chip}{championnat / tournoi}
\mot{gym}{djiim}{salle de sport}
\mot{to work out}{oueu-k AOUT}{faire de l'exercice}
\mot{running / swimming}{RA-n'ning / SOI-m'ning}{course / natation}
\mot{music}{MIou-zik}{musique}
\mot{song / singer / band}{song / SI-ngheu / band}{chanson / chanteur / groupe}
\mot{to sing / dance / play}{sing / daans / plé-i}{chanter / danser / jouer}
\mot{film / movie}{film / MOU-vi}{film (BrE / AmE)}
\mot{cinema / theatre}{SI-neu-ma}{cinéma / théâtre}
\mot{actor / actress}{AK-teu}{acteur / actrice}
\mot{series / episode}{SÍ-eu-ri-z}{série / épisode}
\mot{documentary}{do-kieu-MÈN-t'ri}{documentaire --- \en{watch a documentary}}
\mot{subtitles}{SEUB-taï-tlz}{sous-titres}
\mot{screen / sound / picture}{skriin / saound / PIK-tcheu}{écran / son / image}
\mot{book / novel / story}{bouk / NO-v'l / STOO-ri}{livre / roman / histoire}
\mot{to read / write}{riid / raït}{lire / écrire --- \en{read, read ; wrote, written}}
\mot{to paint / draw}{pé-int / droo}{peindre / dessiner}
\mot{to travel / visit}{TRA-v'l / VI-zit}{voyager / visiter}
\mot{to cook / bake}{kouk / bé-ik}{cuisiner / pâtisser}
\mot{to take photos}{té-ik FO-tôuz}{prendre des photos --- \en{take pictures}}
\end{tabvocab}

\subsection{Thème 16 --- Les émotions et le caractère}
\begin{tabvocab}
\mot{happy / glad}{HA-pi / glad}{heureux, content}
\mot{unhappy / sad}{en-HA-pi / sad}{malheureux, triste}
\mot{angry}{AEN-gri}{en colère --- \en{to be angry with somebody}}
\mot{annoyed}{e-NOÏD}{agacé}
\mot{worried}{OUA-rid}{inquiet}
\mot{afraid / scared}{e-FRÉ-ID / ské-eud}{effrayé --- \en{afraid of}}
\mot{surprised}{seu-PRAÏZD}{surpris}
\mot{excited}{ik-SAÏ-tid}{enthousiaste --- \en{I'm excited about it}}
\mot{bored}{bood}{qui s'ennuie}
\mot{tired / exhausted}{TAÏ-eud / ig-ZOOS-tid}{fatigué / épuisé}
\mot{proud}{praoud}{fier --- \en{proud of}}
\mot{ashamed}{e-CHÉ-IMD}{honteux}
\mot{jealous}{DJÉ-leuss}{jaloux}
\mot{disappointed}{di-ss'e-POÏN-tid}{déçu}
\mot{confident}{KON-fi-d'nt}{confiant}
\mot{nervous}{NEU-veuss}{nerveux, stressé}
\mot{calm}{kaam}{calme}
\mot{lonely}{LÔUN-li}{seul (sentiment)}
\mot{grateful / thankful}{GRÉ-i-tful}{reconnaissant}
\mot{kind / generous}{kaïnd / DJÉ-neu-reus}{gentil / généreux}
\mot{polite / rude}{peu-LAÏT / rood}{poli / impoli}
\mot{honest / dishonest}{O-nist}{honnête / malhonnête}
\mot{hard-working / lazy}{haad-OUOE-king / LÉ-i-zi}{travailleur / paresseux}
\mot{serious / funny}{SÍ-eu-ri-eus / FA-ni}{sérieux / drôle}
\mot{friendly / shy}{FRÈND-li / chaï}{sympathique / timide}
\mot{patient / impatient}{PÉ-i-cheunt}{patient / impatient}
\mot{clever / smart / stupid}{KLÉ-veu / smaart / STiou-pid}{intelligent / futé / stupide}
\mot{brave / cowardly}{bré-iv}{courageux / lâche}
\mot{talkative / quiet}{TOO-keu-tiv / KOUAÏ-et}{bavard / calme}
\mot{ambitious}{am-BI-cheuss}{ambitieux}
\mot{feelings}{FII-lingz}{sentiments}
\mot{mood}{mouud}{humeur --- \en{in a good mood}}
\end{tabvocab}

\subsection{Thème 17 --- Technologie, téléphone, Internet}
\begin{tabvocab}
\mot{phone / mobile / cell}{fôoun / MÔ-baïl}{téléphone / portable}
\mot{screen}{skriin}{écran}
\mot{battery}{BA-teu-ri}{batterie}
\mot{charger}{TCHAAD-jeu}{chargeur}
\mot{to charge / unplug}{tchaadj / eun-PLEUG}{charger / débrancher}
\mot{to call / ring}{kooll / ring}{appeler --- \en{to give someone a call}}
\mot{text message / SMS}{tèkst MÈ-ssidj}{message texte}
\mot{to send / receive}{sènd / ri-SII-uv}{envoyer / recevoir}
\mot{call}{kooll}{appel --- \en{missed call}}
\mot{network / signal}{NÈT-oueu-k / SIG-neul}{réseau / signal}
\mot{data / Wi-Fi}{DÉ-i-ta}{données / wi-fi}
\mot{computer / laptop}{keum-Piou-teu}{ordinateur / portable}
\mot{keyboard / mouse}{KII-bood / maouss}{clavier / souris}
\mot{file / folder}{faïl / FÔUL-deu}{fichier / dossier}
\mot{to save / delete / download}{sé-iv / di-LIIT / DAOUN-lôd}{enregistrer / supprimer / télécharger}
\mot{to log in / out}{log IN / AOUT}{se connecter / se déconnecter}
\mot{password}{PAASS-oueu-d}{mot de passe}
\mot{username}{Iou-zeu-NÉ-im}{nom d'utilisateur}
\mot{website}{OUÈB-saït}{site web}
\mot{search engine}{seutch ÈN-djin}{moteur de recherche}
\mot{app / software / hardware}{ap / SOFT-oué-eu}{application / logiciel / matériel}
\mot{email / inbox}{ii-mé-il / IN-boks}{courriel / boîte de réception}
\mot{attachment}{e-TATCH-m'nt}{pièce jointe}
\mot{social media}{SÔ-cheul MII-di-a}{réseaux sociaux}
\mot{account}{e-KAOUNT}{compte}
\mot{to post / share / like / follow}{pôst / ché-eu}{publier / partager / aimer / suivre}
\mot{to click / tap / scroll}{klik / tap / skroul}{cliquer / taper / faire défiler}
\mot{video call}{VI-di-ou kooll}{appel vidéo}
\mot{to take a screenshot}{ské-in-chot}{faire une capture d'écran}
\mot{to fix / repair}{fiks / ri-PÉ-eu}{réparer}
\end{tabvocab}

\subsection{Thème 18 --- Médias, actualité, société}
\begin{tabvocab}
\mot{news}{niouz}{actualités (singulier !) --- \en{The news is bad}}
\mot{newspaper}{NIouz-pé-i-peu}{journal}
\mot{TV / channel}{téli-vi / TCHA-n'l}{télévision / chaîne}
\mot{radio}{RÉ-i-di-ou}{radio}
\mot{press}{prèss}{presse}
\mot{journalist / reporter}{DJEU-neu-list}{journaliste / reporter}
\mot{headline}{HÈD-laïn}{gros titre}
\mot{article}{AA-ti-kl}{article}
\mot{interview}{IN-teu-viou}{entretien, interview}
\mot{advertisement / advert}{ed-VEU-tiss-m'nt}{publicité}
\mot{information}{in-feu-MÉ-i-cheun}{informations (indénombrable)}
\mot{rumour / fake news}{ROU-meu}{rumeur / fausse nouvelle}
\mot{government}{GOU-veun-m'nt}{gouvernement}
\mot{president / minister}{PRÉ-zi-d'nt}{président / ministre}
\mot{election / vote}{i-LÈK-cheun / vôt}{élection / vote}
\mot{law / rule}{loo / roul}{loi / règle}
\mot{rights}{raïts}{droits --- \en{human rights}}
\mot{freedom}{FRII-deum}{liberté}
\mot{citizen}{SI-ti-z'n}{citoyen}
\mot{society}{seu-SAÏ-eu-ti}{société}
\mot{community}{keu-MIou-ni-ti}{communauté, quartier}
\mot{population}{po-piou-LÉ-i-cheun}{population}
\mot{development}{di-VÉ-leup-m'nt}{développement}
\mot{economy / market}{i-KO-no-mi}{économie / marché}
\mot{poverty / wealth}{PO-veu-ti / ouèlth}{pauvreté / richesse}
\mot{unemployment}{eun-im-PLOÏ-m'nt}{chômage}
\mot{education / health}{é-djou-KÉ-i-cheun / hèlth}{éducation / santé}
\mot{environment}{in-VAÏ-reun-m'nt}{environnement}
\mot{crime / theft}{kraïm / thèft}{crime / vol}
\mot{police / court / prison}{peu-LIISS / koort / PRÍ-z'n}{police / tribunal / prison}
\end{tabvocab}

\subsection{Thème 19 --- Religion et vie spirituelle}
\begin{tabvocab}
\mot{religion}{ri-LI-djeun}{religion}
\mot{God}{god}{Dieu}
\mot{faith}{fé-ith}{foi}
\mot{to believe in}{bi-LII-uv in}{croire en}
\mot{prayer}{pré-eu}{prière --- \en{to say a prayer}}
\mot{to pray}{pré-i}{prier}
\mot{worship}{OUOE-chip}{culte, adoration}
\mot{mosque / church / temple}{mosk / tcheutch / TÈM-pl}{mosquée / église / temple}
\mot{imam / priest / pastor}{i-MAAM / priist / PAASS-teu}{imam / prêtre / pasteur}
\mot{holy book}{HÔ-li bouk}{livre saint}
\mot{verse / chapter}{veuss / TCHA-p-teu}{verset / chapitre (sourate : \en{surah})}
\mot{recitation}{rè-si-TÉ-i-cheun}{récitation}
\mot{pilgrimage}{PI-lgri-midj}{pèlerinage}
\mot{fasting}{FAAS-ting}{jeûne --- \en{to fast}}
\mot{charity / almsgiving}{TCHA-ri-ti}{aumône, charité}
\mot{blessing}{BLÈ-ssing}{bénédiction}
\mot{to bless}{blèss}{bénir}
\mot{soul}{sôoul}{âme}
\mot{sin / forgiveness}{sine / feu-GIV-n's}{péché / pardon}
\mot{to repent / forgive}{ri-PÈNT / feu-GIV}{se repentir / pardonner}
\mot{mercy}{MEU-si}{miséricorde}
\mot{peace}{piiss}{paix --- \en{Peace be upon you}}
\mot{religious / secular}{ri-LI-djeus}{religieux / laïque}
\mot{celebration / festival}{sè-leu-BRÉ-i-cheun}{célébration / fête}
\end{tabvocab}

\begin{astuce}[Formules de respect interculturel]
\en{Peace be upon you.} (Que la paix soit sur vous.) \quad
\en{God willing.} / \en{If God wills it.} (Si Dieu le veut.) \quad
\en{May God bless you.} (Que Dieu vous bénisse.) \quad
\en{Congratulations on your wedding.} (Félicitations pour votre mariage.) \quad
\en{My condolences.} / \en{I'm sorry for your loss.} (Mes condoléances.) \quad
\en{Ramadan Mubarak.} / \en{Happy Tabaski!} / \en{Merry Christmas!}
\end{astuce}

\subsection{Thème 20 --- Décrire, expliquer, nuancer (mots-outils)}
\begin{tabvocab}
\mot{thing}{thing}{chose --- \en{the most important thing}}
\mot{way}{oué-i}{façon, manière --- \en{in this way, by the way}}
\mot{reason}{RII-z'n}{raison --- \en{the reason why}}
\mot{idea}{aï-DÍ-eu}{idée}
\mot{example}{ig-ZAM-pl}{exemple --- \en{for example}}
\mot{fact}{fakt}{fait --- \en{in fact}}
\mot{opinion}{e-PÍ-nieun}{opinion --- \en{in my opinion}}
\mot{point}{poïnt}{point, argument --- \en{good point}}
\mot{meaning}{MII-ning}{sens, signification}
\mot{difference}{DI-feu-reuns}{différence --- \en{make a difference}}
\mot{problem / solution}{PRO-b'l / se-LIOU-cheun}{problème / solution}
\mot{question / answer}{KWÈS-tchoun}{question / réponse}
\mot{result}{ri-ZEULT}{résultat}
\mot{chance}{tchaans}{chance, occasion --- \en{give me a chance}}
\mot{choice}{tchoïs}{choix}
\mot{decision}{di-SÍ-jeun}{décision}
\mot{advice}{ed-VAÏS}{conseil (indénombrable) --- \en{a piece of advice}}
\mot{news}{niouz}{nouvelle (singulier)}
\mot{experience}{ik-SPÍ-eu-ri-euns}{expérience}
\mot{knowledge}{NO-lidj}{connaissances}
\mot{purpose}{PEU-peuss}{but --- \en{on purpose} exprès}
\mot{advantage / disadvantage}{ed-VAAN-tidj}{avantage / inconvénient}
\mot{as soon as possible}{euss SOUN euss}{dès que possible --- \en{ASAP}}
\mot{at least / at most}{at liist}{au moins / au plus}
\mot{for instance}{fo IN-ss'teuns}{par exemple}
\mot{by the way}{baï dhe oué-i}{au fait, à propos}
\mot{in other words}{in EUDheu oueu-dz}{autrement dit}
\mot{that's why}{dhats ouaï}{c'est pourquoi}
\mot{it depends}{it di-PÈNDZ}{cela dépend}
\mot{I mean}{aï miin}{je veux dire}
\mot{actually}{AK-tchou-a-li}{en fait (faux ami : pas « actuellement »)}
\mot{probably / maybe / perhaps}{PRO-beu-bli}{probablement / peut-être / peut-être}
\end{tabvocab}

\section{Les 98 verbes essentiels (à savoir conjuguer par c\oe ur)}
\begin{tableau}[98 verbes essentiels avec leurs trois formes]
\footnotesize
\begin{longtable}{@{}p{2.4cm}p{2.2cm}p{2.4cm}p{3.0cm}p{2.4cm}p{2.2cm}@{}}
\entete{Base} & \entete{Passé} & \entete{Participe} & \entete{Sens} & \entete{Base} & \entete{Passé / p.p.} \tabularnewline
\endfirsthead
\entete{Base} & \entete{Passé} & \entete{Participe} & \entete{Sens} & \entete{Base} & \entete{Passé / p.p.} \tabularnewline
\endhead
\hline
\en{be} & \en{was/were} & \en{been} & être & \en{make} & \en{made / made} \tabularnewline
\en{have} & \en{had} & \en{had} & avoir & \en{take} & \en{took / taken} \tabularnewline
\en{do} & \en{did} & \en{done} & faire & \en{come} & \en{came / come} \tabularnewline
\en{say} & \en{said} & \en{said} & dire & \en{go} & \en{went / gone} \tabularnewline
\en{get} & \en{got} & \en{got(ten)} & obtenir, devenir & \en{see} & \en{saw / seen} \tabularnewline
\en{know} & \en{knew} & \en{known} & savoir & \en{want} & \en{wanted / wanted} \tabularnewline
\en{think} & \en{thought} & \en{thought} & penser & \en{need} & \en{needed / needed} \tabularnewline
\en{tell} & \en{told} & \en{told} & dire à qqn & \en{feel} & \en{felt / felt} \tabularnewline
\en{find} & \en{found} & \en{found} & trouver & \en{try} & \en{tried / tried} \tabularnewline
\en{give} & \en{gave} & \en{given} & donner & \en{leave} & \en{left / left} \tabularnewline
\en{work} & \en{worked} & \en{worked} & travailler & \en{become} & \en{became / become} \tabularnewline
\en{call} & \en{called} & \en{called} & appeler & \en{put} & \en{put / put} \tabularnewline
\en{ask} & \en{asked} & \en{asked} & demander & \en{mean} & \en{meant / meant} \tabularnewline
\en{use} & \en{used} & \en{used} & utiliser & \en{keep} & \en{kept / kept} \tabularnewline
\en{look} & \en{looked} & \en{looked} & regarder & \en{let} & \en{let / let} \tabularnewline
\en{seem} & \en{seemed} & \en{seemed} & sembler & \en{begin} & \en{began / begun} \tabularnewline
\en{help} & \en{helped} & \en{helped} & aider & \en{show} & \en{showed / shown} \tabularnewline
\en{talk} & \en{talked} & \en{talked} & parler & \en{hear} & \en{heard / heard} \tabularnewline
\en{turn} & \en{turned} & \en{turned} & tourner & \en{play} & \en{played / played} \tabularnewline
\en{start} & \en{started} & \en{started} & commencer & \en{run} & \en{ran / run} \tabularnewline
\en{move} & \en{moved} & \en{moved} & déménager, bouger & \en{answer} & \en{answered / answered} \tabularnewline
\en{like} & \en{liked} & \en{liked} & aimer & \en{believe} & \en{believed / believed} \tabularnewline
\en{live} & \en{lived} & \en{lived} & habiter & \en{bring} & \en{brought / brought} \tabularnewline
\en{love} & \en{loved} & \en{loved} & aimer (fort) & \en{happen} & \en{happened / happened} \tabularnewline
\en{eat} & \en{ate} & \en{eaten} & manger & \en{write} & \en{wrote / written} \tabularnewline
\en{drink} & \en{drank} & \en{drunk} & boire & \en{sit} & \en{sat / sat} \tabularnewline
\en{sleep} & \en{slept} & \en{slept} & dormir & \en{stand} & \en{stood / stood} \tabularnewline
\en{read} & \en{read} & \en{read} & lire & \en{lose} & \en{lost / lost} \tabularnewline
\en{buy} & \en{bought} & \en{bought} & acheter & \en{pay} & \en{paid / paid} \tabularnewline
\en{sell} & \en{sold} & \en{sold} & vendre & \en{meet} & \en{met / met} \tabularnewline
\en{open} & \en{opened} & \en{opened} & ouvrir & \en{include} & \en{included / included} \tabularnewline
\en{close} & \en{closed} & \en{closed} & fermer & \en{continue} & \en{continued / continued} \tabularnewline
\en{learn} & \en{learnt / learned} & \en{learnt} & apprendre & \en{set} & \en{set / set} \tabularnewline
\en{teach} & \en{taught} & \en{taught} & enseigner & \en{change} & \en{changed / changed} \tabularnewline
\en{understand} & \en{understood} & \en{understood} & comprendre & \en{lead} & \en{led / led} \tabularnewline
\en{remember} & \en{remembered} & \en{remembered} & se souvenir & \en{explain} & \en{explained / explained} \tabularnewline
\en{forget} & \en{forgot} & \en{forgotten} & oublier & \en{watch} & \en{watched / watched} \tabularnewline
\en{wait} & \en{waited} & \en{waited} & attendre & \en{follow} & \en{followed / followed} \tabularnewline
\en{stop} & \en{stopped} & \en{stopped} & arrêter & \en{finish} & \en{finished / finished} \tabularnewline
\en{walk} & \en{walked} & \en{walked} & marcher & \en{create} & \en{created / created} \tabularnewline
\en{drive} & \en{drove} & \en{driven} & conduire & \en{speak} & \en{spoke / spoken} \tabularnewline
\en{fly} & \en{flew} & \en{flown} & voler & \en{carry} & \en{carried / carried} \tabularnewline
\en{grow} & \en{grew} & \en{grown} & grandir, cultiver & \en{allow} & \en{allowed / allowed} \tabularnewline
\en{build} & \en{built} & \en{built} & construire & \en{add} & \en{added / added} \tabularnewline
\en{send} & \en{sent} & \en{sent} & envoyer & \en{spend} & \en{spent / spent} \tabularnewline
\en{break} & \en{broke} & \en{broken} & casser & \en{cut} & \en{cut / cut} \tabularnewline
\en{catch} & \en{caught} & \en{caught} & attraper & \en{raise} & \en{raised / raised} \tabularnewline
\en{choose} & \en{chose} & \en{chosen} & choisir & \en{kill} & \en{killed / killed} \tabularnewline
\en{cost} & \en{cost} & \en{cost} & coûter & \en{die} & \en{died / died} \tabularnewline
\end{longtable}
\end{tableau}

\begin{mission}[Chapitre 12 --- votre premier millier]
Objectif chiffré : \textbf{1\,000 mots en 100 jours}. Procédure :\\
\textbf{1.} Choisissez un thème par jour (20 thèmes ci-dessus, puis répétez).\\
\textbf{2.} Sélectionnez 15 mots du thème --- ceux dont vous avez vraiment besoin.\\
\textbf{3.} Créez 15 cartes Anki, chacune avec \textbf{une phrase personnelle} et la
prononciation.\\
\textbf{4.} Dites les 15 phrases à voix haute, deux fois.\\
\textbf{5.} Le soir, écrivez 5 phrases de journal en utilisant au moins 5 mots du jour.\\
En 100 jours, vous aurez 1\,000 mots actifs --- c'est-à-dire 85\,\% de la conversation
quotidienne.
\end{mission}
